\section{分离公理}\label{section:ch04-separation-axioms}

本节将引进三个分离公理，并阐明它们的一些性质。前面已经介绍过 Hausdorff 公理，另一些公理与它类似，只是比它更强一些。每当引进新概念之后，我们总是讨论它们与本书已经给出的那些公理和概念之间的关系。

我们讲过，空间 $X$ 称为 Hausdorff 的，指的是：如果对于 $X$ 中每两个互不相同的点 $x$ 和 $y$，存在无交的两个开集分别包含 $x$ 和 $y$。

\begin{definition}\label{definition:ch04-regular-normal}
设 $X$ 中的每一个单点集在 $X$ 中都是闭的。如果对于任意给定的一个点 $x$ 和不包含这个点的一个闭集 $B$，存在无交的两个开集分别包含 $x$ 和 $B$，则称 $X$ 为\textbf{正则的}（regular）。如果对于 $X$ 中每一对无交的闭集 $A$ 和 $B$ 总存在无交的开集分别包含它们，则称 $X$ 是\textbf{正规的}（normal）。
\end{definition}

显然，正则空间是 Hausdorff 的，正规空间是正则的。（为此必须将单点集是闭的这个条件作为正则性和正规性定义的一部分，具有平庸拓扑的两点集空间，虽然满足正则性和正规性定义中的其余条件，但它却不是 Hausdorff 的。）下面的例 1 和例 3 表明了正则性公理比 Hausdorff 公理更强，正规性公理比正则性公理更强。

这些公理之所以被称为分离性公理，是因为它们都涉及用无交的开集，把一定类型的集合彼此``分离''。当然，在学习连通空间之前，我们已经使用过``分割''这个词。但在那里，我们也是试图找到一些无交的开集，使得它们的并就是整个空间。现在所讨论的与此不同，因为这里所说的开集并不要求满足上述条件。

这三个分离公理如\cref{figure:ch04-separation-axioms}所示。

\par
\begin{figure}[htbp]
\centering
\includegraphics[width=0.7\textwidth]{ch04-fig31-1.png}
\caption{Hausdorff、正则与正规}\label{figure:ch04-separation-axioms}
\end{figure}

还有其他的方式刻画分离性公理。其中常用的一种如下述引理所述：

\begin{lemma}\label{lemma:ch04-regular-neighborhood}
设 $X$ 是一个拓扑空间，$X$ 中的单点集都是闭的。则
\begin{enumerate}
\item $X$ 是正则的当且仅当对 $X$ 中任意给定的一个点 $x$ 和 $x$ 的任何一个邻域 $U$，存在 $x$ 的一个邻域 $V$，使得 $\overline{V}\subset U$.
\item $X$ 是正规的当且仅当对于任意闭集 $A$ 和包含 $A$ 的任何一个开集 $U$，存在一个包含 $A$ 的开集 $V$，使得 $\overline{V}\subset U$.
\end{enumerate}
\end{lemma}

\begin{proof}
(a) 设 $X$ 是正则的，给定点 $x$ 和 $x$ 的邻域 $U$。令 $B=X\setminus U$，则 $B$ 是一个闭集。根据假设，存在分别包含 $x$ 和 $B$ 的无交开集 $V$ 和 $W$。因为若 $y\in B$，则集合 $W$ 就是 $y$ 的一个邻域，它与 $V$ 无交，所以集合 $\overline{V}$ 与 $B$ 无交。因此 $\overline{V}\subset U$。

为了证明其逆命题成立，设给定点 $x$ 和不包含 $x$ 的闭集 $B$。令 $U=X\setminus B$。根据假设，存在 $x$ 的邻域 $V$，使得 $\overline{V}\subset U$。于是，开集 $V$ 和 $X\setminus\overline{V}$ 就是分别包含 $x$ 和 $B$ 的无交开集。从而 $X$ 是正则的。

(b) 这个证明完全与上面的证明相同，只需在整个证明中用集合 $A$ 代替点 $x$ 就可以了。
\end{proof}

现在来研究分离性公理与原先引进的一些概念之间的联系。

\begin{theorem}\label{theorem:ch04-hausdorff-regular-hereditary}
\begin{enumerate}
\item Hausdorff 空间的子空间是 Hausdorff 的。Hausdorff 空间的积空间也是 Hausdorff 的.
\item 正则空间的子空间是正则的。正则空间的积空间也是正则的.
\end{enumerate}
\end{theorem}

\begin{proof}
(a) 这是第 17 节习题中的结论，在这里我们给出证明。设 $X$ 是一个 Hausdorff 空间，$x$ 和 $y$ 是 $X$ 的子空间 $Y$ 中的两点。如果 $U$ 和 $V$ 分别是点 $x$ 和 $y$ 在 $X$ 中的无交邻域，那么 $U\cap Y$ 和 $V\cap Y$ 便分别是 $x$ 和 $y$ 在 $Y$ 中的无交邻域。

设 $\{X_{\alpha}\}$ 是 Hausdorff 的空间的一个族。设 $\boldsymbol{x}=(x_{\alpha})$ 和 $\boldsymbol{y}=(y_{\alpha})$ 是积空间 $\prod X_{\alpha}$ 中不同的两点。因为 $\boldsymbol{x}\neq\boldsymbol{y}$，故存在某一个指标 $\beta$，使得 $x_{\beta}\neq y_{\beta}$。在 $X_{\beta}$ 中选取分别包含 $x_{\beta}$ 和 $y_{\beta}$ 的无交开集 $U$ 和 $V$。这样，集合 $\pi_{\beta}^{-1}(U)$ 和 $\pi_{\beta}^{-1}(V)$ 就是 $\prod X_{\alpha}$ 中分别包含 $\boldsymbol{x}$ 和 $\boldsymbol{y}$ 的无交开集。

(b) 设 $Y$ 是一个正则空间 $X$ 的子空间。则 $Y$ 中的单点集都是闭的。设 $x$ 是 $Y$ 的一个点，$B$ 是 $Y$ 中不包含 $x$ 的一个闭子集。于是 $\overline{B}\cap Y=B$，其中 $\overline{B}$ 表示 $B$ 在 $X$ 中的闭包。因此，$x\notin\overline{B}$。再应用 $X$ 的正则性，我们可以选取 $X$ 中分别包含 $x$ 和 $\overline{B}$ 的无交开集 $U$ 和 $V$。因此 $U\cap Y$ 和 $V\cap Y$ 分别是 $Y$ 中包含 $x$ 和 $B$ 的无交开集。

设 $\{X_{\alpha}\}$ 是正则空间的一个族。令 $X=\prod X_{\alpha}$。根据 (a) 可见，$X$ 是一个 Hausdorff 空间，因此 $X$ 中的单点集都是闭的。我们应用前面的引理来证明 $X$ 的正则性。设 $\boldsymbol{x}=(x_{\alpha})$ 为 $X$ 的一个点，$U$ 为 $\boldsymbol{x}$ 在 $X$ 中的一个邻域。取基中的元素 $\prod U_{\alpha}$ 使得它包含于 $U$ 中，且包含着点 $\boldsymbol{x}$。对于每一个 $\alpha$，选取 $x_{\alpha}$ 在 $X_{\alpha}$ 中的一个邻域 $V_{\alpha}$，使得 $\overline{V_{\alpha}}\subset U_{\alpha}$。如果恰好有 $U_{\alpha}=X_{\alpha}$，那就选取 $V_{\alpha}=X_{\alpha}$。于是 $V=\prod V_{\alpha}$ 便是点 $\boldsymbol{x}$ 在 $X$ 中的一个邻域。根据乘积的闭包等于各因子闭包的乘积这一事实，$\overline{V}=\prod\overline{V_{\alpha}}$，所以 $\overline{V}\subset\prod U_{\alpha}\subset U$。因此，$X$ 是正则的。
\end{proof}

稍后我们将在本节和下一节中看到，对于正规空间却没有相似的定理。

\begin{example}\label{example:ch04-rk-not-regular}
空间 $\mathbb{R}_{K}$ 是 Hausdorff 的，但不是正则的。前面讲过 $\mathbb{R}_{K}$ 表示实直线，它以所有开区间 $(a,b)$ 和形如 $(a,b)-K$ 的集合为其拓扑的一个基，其中 $K=\{\,1/n\mid n\in\mathbb{Z}_{+}\,\}$。这个空间是 Hausdorff 的，因为任意两个不同的点都存在两个无交的开区间包含它们。

但这个空间不是正则的。集合 $K$ 是 $\mathbb{R}_{K}$ 的一个闭集，并且不包含 $0$ 点。设存在无交的开集 $U$ 和 $V$ 分别包含 $0$ 和 $K$。选取基中包含于 $U$ 的一个元素，使得它包含着 $0$，那么基中的这个元素必定形如 $(a,b)-K$，因为基中每一个形如 $(a,b)$ 并包含 $0$ 的元素都与 $K$ 相交。选取足够大的 $n$ 使得 $1/n\in(a,b)$，再取基中包含于 $V$ 且包含 $1/n$ 的一个元素，那么基中这个元素必形如 $(c,d)$。最后，取点 $z$，使得 $z<1/n$，且 $z>\max\{c,1/(n+1)\}$。这样 $z$ 既属于 $U$ 又属于 $V$，因此这两个集合并非无交。如\cref{figure:ch04-rk-not-regular}所示。
\end{example}

\par
\begin{figure}[htbp]
\centering
\includegraphics[width=0.6\textwidth]{ch04-fig31-2.png}
\caption{$\mathbb{R}_{K}$ 不是正则的}\label{figure:ch04-rk-not-regular}
\end{figure}

\begin{example}\label{example:ch04-sorgenfrey-line-normal}
空间 $\mathbb{R}_{\ell}$ 是正规的。易见单点集都是 $\mathbb{R}_{\ell}$ 中的闭集，这是由于 $\mathbb{R}_{\ell}$ 的拓扑比 $\mathbb{R}$ 的拓扑更细的缘故。为了证明正规性，设 $A$ 和 $B$ 是 $\mathbb{R}_{\ell}$ 中无交的两个闭集。对 $A$ 中的每一个点 $a$ 选取一个基元素 $[a,x_{a})$，使得它与 $B$ 无交。对于 $B$ 中的每一个点 $b$，选取一个基元素 $[b,x_{b})$，使得它与 $A$ 无交。于是开集
\begin{align*}
U=\bigcup_{a\in A}[a,x_{a})\quad\text{和}\quad V=\bigcup_{b\in B}[b,x_{b})
\end{align*}
分别包含集合 $A$ 和 $B$，并且是无交的。
\end{example}

\begin{example}\label{example:ch04-sorgenfrey-plane-not-normal}
Sorgenfrey 平面 $\mathbb{R}_{\ell}^{2}$ 不是正规的。

空间 $\mathbb{R}_{\ell}$ 是正则的（实际上是正规的），因此积空间 $\mathbb{R}_{\ell}^{2}$ 是正则的。于是，给出这个例子可以达到两个目的。一是说明正则空间未必是正规的，二是说明两个正规空间的积空间未必是正规的。

我们假设 $\mathbb{R}_{\ell}^{2}$ 是正规的，并且由此推出矛盾。设 $L$ 是由 $\mathbb{R}_{\ell}^{2}$ 中所有形如 $x\times(-x)$ 的点构成的一个子空间。$L$ 在 $\mathbb{R}_{\ell}^{2}$ 中是闭的，并且具有离散拓扑。因 $L$ 中的每一个子集 $A$ 在 $L$ 中都是闭的，所以在 $\mathbb{R}_{\ell}^{2}$ 中也都是闭的。又因为 $L-A$ 在 $\mathbb{R}_{\ell}^{2}$ 中也是闭的，所以对于 $L$ 的任意非空真子集 $A$ 都能找到分别包含 $A$ 和 $L-A$ 的无交开集 $U_{A}$ 和 $V_{A}$。

设 $D$ 是 $\mathbb{R}_{\ell}^{2}$ 中坐标为有理数的点构成的集合。$D$ 在 $\mathbb{R}_{\ell}^{2}$ 中稠密。定义映射 $\theta\colon\mathcal{P}(L)\to\mathcal{P}(D)$ 为：
\begin{align*}
\theta(A)=D\cap U_{A}\quad\text{如果}\quad\varnothing\subsetneq A\subsetneq L,
\end{align*}
\begin{align*}
\theta(\varnothing)=\varnothing,\qquad\theta(L)=D.
\end{align*}
我们将证明 $\theta\colon\mathcal{P}(L)\to\mathcal{P}(D)$ 是一个单射。

设 $A$ 是 $L$ 的一个非空真子集。$\theta(A)=D\cap U_{A}$ 既不是空集（因为 $U_{A}$ 是开集，$D$ 在 $\mathbb{R}_{\ell}^{2}$ 中稠密），也不是 $D$（因为 $D\cap V_{A}\neq\varnothing$）。剩下要证明的是：若 $B$ 是 $L$ 的另外一个非空真子集，则 $\theta(A)\neq\theta(B)$。

存在一点只属于 $A$ 和 $B$ 其中之一而不属于另一个，不妨设 $x\in A$，但 $x\notin B$。于是 $x\in L-B$，因此 $x\in U_{A}\cap V_{B}$。又因这个集合是开的并且非空，所以其中必包含 $D$ 的点。这些点属于 $U_{A}$，却不属于集合 $U_{B}$。因此，$D\cap U_{A}\neq D\cap U_{B}$，即 $\theta$ 是一个单射。

现在我们证明存在单射 $\phi\colon\mathcal{P}(D)\to L$。因为 $D$ 是可数的无限集，$L$ 与 $\mathbb{R}$ 有相同的基数，因此可以定义从 $\mathcal{P}(\mathbb{Z}_{+})$ 到 $\mathbb{R}$ 的一个单射 $\psi$。$\psi$ 的定义如下：将 $\mathbb{Z}_{+}$ 的子集 $S$ 映到无限的十进小数 $.a_{1}a_{2}\cdots$，即
\begin{align*}
\psi(S)=\sum_{i=1}^{\infty}a_{i}/10^{i},
\end{align*}
其中，当 $i\in S$ 时，$a_{i}=0$，当 $i\notin S$ 时，$a_{i}=1$。这时，复合映射
\begin{align*}
\mathcal{P}(L)\xrightarrow{\theta}\mathcal{P}(D)\xrightarrow{\psi}L
\end{align*}
便是从 $\mathcal{P}(L)$ 到 $L$ 的一个单射。根据\cref{theorem:ch01-Cantor} 可见，这个映射是不存在的。因此得到一个矛盾。

证明 $\mathbb{R}_{\ell}^{2}$ 不是正规的所用的这个方法，在某些地方并不是那么让人满意。我们只说明了存在 $L$ 的非空真子集 $A$ 使得集合 $A$ 和 $B=L-A$ 不能包含在 $\mathbb{R}_{\ell}^{2}$ 的两个无交的开集中，却没能说明到底集合 $A$ 是怎样的。事实上，$L$ 中所有坐标为有理数的点构成的集合就是我们要找到那个集合，但其证明却不那么容易。我们将其留作练习。
\end{example}

\subsection*{习题}

\begin{enumerate}
\item 证明：若 $X$ 是正则的，则 $X$ 的任意两点各自有一个邻域，其闭包无交。
\item 证明：若 $X$ 是正规的，则任意两个无交的闭集各自有一个邻域，其闭包无交。
\item 证明：每一个序拓扑都是正则的。
\item 设 $X$ 和 $X'$ 分别表示具有拓扑 $\mathcal{T}$ 和 $\mathcal{T}'$ 的同一个集合，并且 $\mathcal{T}'\supset\mathcal{T}$。如果其中一个是 Hausdorff（或正则，或正规）空间，那么另一个会怎样呢？
\item 设 $f,g\colon X\to Y$ 是连续的，$Y$ 是一个 Hausdorff 空间。证明 $\{\,x\mid f(x)=g(x)\,\}$ 是 $X$ 中的闭集。
\item 设 $p\colon X\to Y$ 是满的连续闭映射。证明：若 $X$ 是正规的，则 $Y$ 也是正规的。[提示：若 $U$ 是包含 $p^{-1}(\{y\})$ 的开集，证明：存在 $y$ 的邻域 $W$，使得 $p^{-1}(W)\subset U$。]
\item 设 $p\colon X\to Y$ 是一个满的连续闭映射，对于每一个 $y\in Y$，$p^{-1}(\{y\})$ 是紧致的。（这种映射称为\textbf{完备映射}。）
\begin{enumerate}
\item 证明：若 $X$ 是 Hausdorff 的，则 $Y$ 也是 Hausdorff 的.
\item 证明：若 $X$ 是正则的，则 $Y$ 也是正则的.
\item 证明：若 $X$ 是局部紧致的，则 $Y$ 也是局部紧致的.
\item 证明：若 $X$ 满足第二可数性公理，则 $Y$ 也满足第二可数性公理。[提示：设 $\mathcal{B}$ 是 $X$ 的一个可数基。对于 $\mathcal{B}$ 的每一个有限子集 $J$，令 $U_{J}$ 为所有形如 $p^{-1}(W)$ 的集合的并，其中 $W$ 是 $Y$ 中的开集，并且 $p^{-1}(W)$ 包含于 $J$ 的元素的并中。]
\end{enumerate}
\item 设 $X$ 是一个空间，$G$ 是一个拓扑群。$G$ 在 $X$ 上的一个\textbf{作用}（action）是指一个连续映射 $\alpha\colon G\times X\to X$，满足条件：
\begin{enumerate}
\item $e\circ x=x$，对于所有 $x\in X$.
\item $g_{1}\circ(g_{2}\circ x)=(g_{1}\circ g_{2})\circ x$，对于所有 $x\in X$ 和 $g_{1},g_{2}\in G$.
\end{enumerate}
其中 $g\circ x$ 表示 $\alpha(g\times x)$。对于所有的 $x$ 和 $g$，定义 $x\sim g\circ x$。所得到的商空间记为 $X/G$，并且称之为作用 $\alpha$ 的\textbf{轨道空间}（orbit space）。

\textbf{定理}\quad 设 $G$ 是一个紧致拓扑群，$X$ 是一个拓扑空间。$\alpha$ 是 $G$ 在 $X$ 上的一个作用。若对于 Hausdorff 条件、正则性、正规性、局部紧致性，以及第二可数性公理这些性质，$X$ 满足其中的一个性质，则 $X/G$ 也满足同样的性质。[提示：见第 26 节的习题 13。]
\item 令 $A$ 为 $\mathbb{R}_{\ell}^{2}$ 中所有形如 $x\times(-x)$ 的点的集合，其中 $x$ 是有理数。令 $B$ 为 $\mathbb{R}_{\ell}^{2}$ 中所有形如 $x\times(-x)$ 的点的集合，其中 $x$ 是无理数。如果 $V$ 是 $\mathbb{R}_{\ell}^{2}$ 中包含 $B$ 的一个开集，证明：不存在包含 $A$ 的开集 $U$ 与 $V$ 无交。证明的过程如下：
\begin{enumerate}
\item 设 $K_{n}$ 为 $[0,1]$ 中所有使得 $[x,x+1/n)\times[-x,-x+1/n)$ 包含于 $V$ 的无理数的集合。证明 $[0,1]$ 是这些集合 $K_{n}$ 和可数多个单点集的并.
\item 应用第 27 节中的习题 5 证明某一个集合 $\overline{K}_{n}$ 包含着 $\mathbb{R}$ 的一个开区间 $(a,b)$.
\item 证明 $V$ 包含着由所有形如 $x\times(-x+\varepsilon)$ 的点 $x$ 构成的开平行四边形，其中 $a<x<b$，$0<\varepsilon<1/n$.
\item 证明：若 $q$ 是一个有理数，$a<q<b$，则 $\mathbb{R}_{\ell}^{2}$ 中的点 $q\times(-q)$ 是 $V$ 的一个极限点.
\end{enumerate}
\end{enumerate}